\begin{frame}{왜 ``시연''이 아니라 ``비교''인가 — 두 가지 이유}
  \begin{columns}[T]
    \begin{column}{0.5\textwidth}
      \kb{첫째, 시연 자체가 어려운 문제가 있다}
      \begin{itemize}
        \item ``더 자연스러운 걸음걸이'', ``더 안전한 주행'',
          ``더 정중한 말투''
        \item \textit{최적} 하나를 정밀하게 만드는 게 아니라
          \textit{주어진 두 가지 중 나은 쪽}을 가리는 데 인간의
          직감이 훨씬 잘 맞음.
        \item ``비교는 쉬우나 절대 정의는 어렵다'' — 연속적인 질을
          이산 판정(2택1)으로 줄이면 사람이 \textbf{일관되게}
          라벨을 댈 수 있음.
      \end{itemize}
    \end{column}
    \begin{column}{0.5\textwidth}
      \kb{둘째(더 근본적): 시연은 상한이 사람}
      \begin{itemize}
        \item BC는 전문가가 보여주지 않은 전략을 \textbf{절대}
          만들어낼 수 없다.
        \item 선호 기반 보상모델은 ``더 나은 쪽''만 알리면,
          \textbf{사람이 시연한 적 없는 새 시도}에도 점수를 줄 수
          있다.
        \item PPO로 최적화하면 사람이 보여준 적 없는 행동 패턴을
          \textbf{발견}할 수 있다.
      \end{itemize}
    \end{column}
  \end{columns}
\end{frame}
